\chapter{Evaluation}\label{ch:eval}

Die Evaluation legt zuerst die Methodik offen und wertet darauf aufbauend die Messreihen aus. Ihre
Leitfrage ist die Cacheline-Awareness: Welche Entwurfsentscheidung verändert die cache-nahen
Messgrößen --- L1-, L2- und Last-Level-Misses, dTLB-Misses, Branch-Misses und die Latenz je
Genus-Interface-Aufruf --- auf welcher Plattform in welchem Maß? Drei Evaluationsobjekte ordnen das
Kapitel: der \emph{Entwurfsraum} (die Hypothesen H1--H3 und die Achsen-Sensitivitätsanalyse), der
\emph{Apparat} selbst (der heuristische Gegenbeweis und der Messfehler, den die Mess-Instanzen
eintragen) und die \emph{Automatisierung} der Evaluation aus der Experiment-XML, denn die Arbeit
automatisiert im Kern genau diese Evaluation. Der Methodik-Teil (Abschnitt~\ref{sec:eval-method})
entfaltet in acht Unterabschnitten Hypothesen und Messgrößen, die XML-getriebene Kette, Workloads und
Datensätze, Plattformen und Regimes, Mess-Ebenen und Zähler, Bau-Umfang, Fairness sowie den
Gegenbeweis; der Ergebnis-Teil (Abschnitt~\ref{sec:eval-results}) führt in sieben Unterabschnitten
die Identität der Ergebnisse, den Datenfluss, die Messreihen, die vorliegenden Messungen, die
Sensitivitätsanalyse, die Gültigkeitsprüfung und die Zusammenfassung.

%% OWNER-VORLAGE: OV-K5-05 --- Kapitel-Aufbau nach der Methodenliteratur (drei Evaluationsobjekte, acht plus
%% sieben Unterabschnitte, Gültigkeitsprüfung, Zusammenfassung); Bestands-Labels unverändert. Widerspruch genügt.
Die Evaluation folgt der in der empirischen Software- und Systemforschung üblichen Kette aus
Forschungsfrage, Hypothese, Messgröße, Versuchsaufbau, Durchführung, Ergebnis, Gültigkeitsprüfung
und Diskussion~\cite{wohlin2012experimentation,shaw2003writing}: Jede Hypothese benennt ihre
Messgröße und ihr Annahmekriterium vor der Messung, jeder Versuchsaufbau seine Plattform, sein Regime
und seine Wiederholungen, und jedes Ergebnis seine
Unsicherheit~\cite{georges2007rigorous,kalibera2013rigorous}. Weil der Gegenstand ein Artefakt ist ---
der Mess-Apparat mit seinen Achsen ---, tritt die Bewertung des Artefakts nach den Regeln der
Design-Science-Forschung\footnote{Design Science: die Forschungsrichtung der Informationssysteme, die
ein konstruiertes Artefakt (hier den Mess-Apparat) selbst zum Gegenstand der Bewertung macht und dafür
Nützlichkeit, Strenge des Vorgehens und Nachvollziehbarkeit verlangt~\cite{hevner2004designscience}.}
hinzu: Nützlichkeit, Strenge und
Nachvollziehbarkeit~\cite{hevner2004designscience,peffers2007dsrm,wieringa2014designscience}.
Leistungsmessungen werden nach den Berichts-Regeln von Hoefler und Belli~\cite{hoefler2015benchmarking}
und Blackburn et~al.~\cite{blackburn2016truth} ausgewiesen: Nenner, Wiederholungen, Streuung,
Plattform und nicht erhobene Werte (n/a) stehen an jedem Ergebnis. Die drei Evaluationsobjekte ordnen
die Unterabschnitte; die Cacheline-Awareness ist ihre gemeinsame Leitachse.

Die Leitachse der Evaluation ist die Cacheline-Awareness: Verglichen werden die cache-nahen
Messgrößen --- L1-, L2- und Last-Level-Misses, dTLB-Misses, Branch-Misses, Energie sowie die
Latenz-Perzentile je Genus-Interface-Funktion; sie sind sechzehn registrierte Mess-Kategorien mit einem
Regime-Ordinal je Kategorie --- je Achsen-Belegung, je Gattungs- und Genus-Interface-Aufruf (die
Kanal-Adresse lautet stets Achse $\to$ Genus $\to$ Kategorie) und je Plattform, also je registrierter
Maschinen-Signatur: prod1 (AMD Zen~5), prod2 (Intel Core i9-12900K, Alder Lake) und das ODROID-Board
mit Intel-Gracemont-Kernen. Die Frage lautet nicht, welcher Algorithmus schneller ist, sondern welche
Achse die cacheline-bezogenen Kosten auf welcher Mikroarchitektur bewegt. Die PMC-Kategorien
(Performance Monitoring Counter, die Hardware-Zähler der Prozessoren) sind dafür die primären
Messgrößen, die Zeit-Kategorien die Gegenprobe; weil Hardware-Zähler nur nach einer Kreuzprobe am
Objekt vertrauenswürdig sind~\cite{weaver2008counters}, legt Abschnitt~\ref{sec:eval-counters} die
Semantik jedes Zählers je Modell offen.

\section{Methodik}\label{sec:eval-method}
\subsection{Hypothesen und Messgrößen}\label{sec:eval-hypotheses}
Drei \emph{Hypothesen} leiten die Auswertung; jede nennt ihre Messgröße und ihr Annahmekriterium vor
der Messung:
\begin{description}
  \item[H1 (Page-Type-Kosten)] Die mittleren Zugriffskosten je Page-Type-Variante unterscheiden sich
  systematisch nach Workload: dichte Page-Typen gewinnen erwartungsgemäß unter Zipf-Verteilung (YCSB-A),
  bereichsoptimierte Pages unter Bereichs-Scans (YCSB-E). Messgrößen sind die Latenz-Perzentile und die
  Cache-Miss-Kategorien je Workload; H1 gilt als gestützt, wenn die Rangfolge der Page-Type-Varianten
  zwischen YCSB-A und YCSB-E auf derselben Plattform kippt und der Abstand das Unsicherheitsband der
  Bau-Varianten (Abschnitt~\ref{sec:eval-counterproof}) übersteigt.
  %% OWNER-VORLAGE: OV-03 --- Annahme-/Ablehnkriterium für H2 (zwei Instrumente: Hand-Audit über sieben
  %% Achsen und tool-berechneter cppcheck-Score in der Mess-Zeile). Default: Rangkorrelation des
  %% tool-berechneten Scores gegen den Durchsatz derselben Zelle, n/a-Profile ausgeschlossen, Hand-Audit
  %% als qualitative Einordnung, Schwelle vor der Messung zu nennen. Echter Entscheid des Owners.
  \item[H2 (Code-Qualität pro Quelle)] Der vom Original-Quell-Repository geerbte Code-Qualitäts-Score
  korreliert messbar mit dem erreichten Durchsatz. Zwei Instrumente tragen den Score: erstens das
  dokumentierte Hand-Audit je Quell-Repository über sieben Bewertungs-Achsen (Stil, Tests,
  Dokumentation, Wartbarkeit, Pflege-Status, Lizenz, Build-System; Skala 1--5) für 22 bewertete
  Quellen --- eine Akte unter den Audits des Cache-Engine-Repositories ---, zweitens das
  tool-berechnete Instrument, das real in die Mess-Zeile eingeht: die gewichtete
  \texttt{cppcheck}-Befunddichte je tausend physischer Zeilen (Gewichte error~3, warning~2,
  performance~1, portability~1, style~0,5) über die vendorten Original-Quellen; die Akte führt 33
  Einträge, davon 22 als n/a, weil die Rekonstruktion keine vendorte Original-Quelle besitzt, und der
  Wert reist als Spalte \texttt{h2\_score} in jeder Mess-Zeile einer SOTA-Reihe mit. Ergänzt wird
  beides um die dokumentierte \texttt{is\_original}-Markierung je Funktion (Anhang~\ref{app:blocks}).
  Annahmekriterium: H2 gilt als gestützt, wenn der tool-berechnete Score über die belegten SOTA-Profile
  mit dem Durchsatz in derselben Zelle (Plattform, Workload, Regime) einen Rangkorrelations-Koeffizienten
  jenseits einer vor der Messung genannten Schwelle zeigt; n/a-Profile gehen nicht in die Korrelation
  ein, und das Hand-Audit ist die qualitative Einordnung.
  %% OWNER-VORLAGE: OV-K5-03 --- H3-Nenner: drei ValueHandle-Varianten (Inline, External, ChainRef) als
  %% Ausschnitt der fünf Registry-Bausteine der Achse value_handle. Default gesetzt; Widerspruch genügt.
  \item[H3 (Inline- vs.\ External- vs.\ Chain-Verteilung)] Die beobachtete Verteilung der drei
  ValueHandle-Varianten korreliert mit der Page-Dichte: dichte Pages bevorzugen Inline, spärliche Pages
  External und PRT-ART-Heuristiken ChainRef. Die Registry der Achse \texttt{value\_handle} führt fünf
  Bausteine (Inline, External-Pool, Immutable-Shared-Ref, Versioned-Pointer, Chain-Ref; nur Inline ist
  golden-verdrahtet); H3 betrachtet die drei für PRT-ART tragenden Varianten als begründeten
  Ausschnitt. Messgröße ist der Anteil je Variante über der Page-Dichte; H3 gilt als gestützt, wenn
  der Anteil monoton mit der Dichte wandert.
\end{description}
Diese \emph{Evaluations-Hypothesen} sind nicht zu verwechseln mit den gleichnamigen
PRT-ART-Telemetrie-Metriken H1--H3 (Abschnitt~\ref{sec:prtart-demo}): Erstere sind prüfbare Aussagen über
den Entwurfsraum, Letztere Mess-Organe des Prüflings, die unter anderem die Daten für diese Prüfung
liefern.

\subsection{Vollautomatische Experiment-Evaluation aus der XML}\label{sec:eval-xml}
Die Evaluation ist selbst Gegenstand der Automatisierung (Abbildung~\ref{fig:eval-xml-kette}): Eine
Anwender-XML mit der Wurzel \texttt{<comdare\_experiment>} deklariert, \emph{was} gemessen wird,
\emph{wo} und \emph{wann} --- nie das \emph{Wie}. Für die Evaluation tragen die Sektionen
\texttt{workloads}, \texttt{datasets}, \texttt{measurement\_tooling}, \texttt{measurement\_categories},
\texttt{run\_methodology}, \texttt{writeback\_methods} und \texttt{output}; der Profil-Validator prüft
sie vor jedem Lauf gegen die Registries (Achsen-Namen und -Werte, Verbund-Bindungen, Fairness-Modi,
Akten). Aus der XML entsteht deterministisch der Plan --- der Director-Walk des Planers, der als reine
Text-Emission auch ohne Bau abrufbar ist ---, daraus die Bau-Jobs des CacheEngineBuilders und die
Tier-Binaries, deren Messung in der Mappe (xlsx) landet; die CSV ist die Projektion der Mappe, und die
Generatoren des Werkzeug-Repositories (Binary-zu-CSV, CSV-zu-LaTeX, Diagramm-Generator,
Anhang-Generator) rendern Steckbriefe, Tabellen und TikZ-Diagramme je Sprache in ein
\texttt{tabellen}-Verzeichnis, das das Manuskript einbindet. Kein Schritt dieser Kette wird von Hand
ausgeführt; jede Zahl im Anhang trägt den Lauf, der sie erzeugt hat (Anhang~\ref{app:measurements}).
Durchführungs-Stand: Die Kette ist am Durchstich real gelaufen --- die Smoke-Messreihe über 43
Permutationen im Anhang ist ihr Erzeugnis ---; die XML-Sektion \texttt{<publish>} für den
Lager-Rückschrieb, der kommandozeilengeführte Erstellungsmodus der XML
(Abschnitt~\ref{sec:impl-pipeline}) und die dritte Ausgabe-Factory LaTeX
(Abschnitt~\ref{sec:eval-dataflow}) sind Ausbaustufen, nicht Behauptungen.
\begin{figure}[htbp]\centering
% Kein \resizebox: fuenf Kacheln je 2.55cm Textbreite im Raster 3.1cm (Gesamtbreite ~15.1cm) passen nativ
% in \textwidth; Regel: resizebox nur zum VERKLEINERN (vgl. fig:m-model).
\begin{tikzpicture}[font=\footnotesize,>=Stealth,
  st/.style={draw,rounded corners,fill=blue!7,text width=2.55cm,minimum height=1.1cm,align=center,inner sep=2pt,font=\scriptsize},
  og/.style={st,fill=green!9},
  pl/.style={st,fill=orange!10,dashed,minimum height=0.8cm,font=\tiny}]
  \node[st] (x) at (0,0){\textbf{Anwender-XML}\\\texttt{<comdare\_\allowbreak{}experiment>}: workloads, datasets, tooling, categories, run\_methodology, writeback, output};
  \node[st] (v) at (3.1,0){\textbf{Validator}\\Registries, Akten, Verbund-Bindung, Fairness-Modi};
  \node[st] (p) at (6.2,0){\textbf{Planer}\\Director-Walk, Plan-Text, Bau-Legende};
  \node[st] (c) at (9.3,0){\textbf{CEB}\\Bau-Jobs je System-Permutation, Tier-Binaries};
  \node[st] (m) at (12.4,0){\textbf{Messung}\\wallclock / macro / micro, PMC-Zähler, ein Mess-Thread};
  \node[og] (xl) at (12.4,-2.6){\textbf{Mappe (xlsx)}\\Stamm: entsteht immer; ein Datenblatt (konstanter Blattschlüssel); Blatt je Unter-Achsen-Permutation: Soll};
  \node[og] (cs) at (9.3,-2.6){\textbf{Projektion (csv)}\\Kind der Mappe; WIDE 189 Spalten, Pipeline-Sicht 16 + 16};
  \node[og] (g) at (6.2,-2.6){\textbf{Generatoren}\\Binary-zu-CSV, CSV-zu-LaTeX, Diagramme, Anhang};
  \node[og] (t) at (3.1,-2.6){\textbf{tabellen/}\\je Sprache: Steckbriefe, booktabs-Tabellen, TikZ};
  \node[og] (ms) at (0,-2.6){\textbf{Manuskript}\\\textbackslash{}input im Anhang; jede Zahl mit ihrem Lauf};
  \foreach \a/\b in {x/v,v/p,p/c,c/m}{\draw[->,thick] (\a)--(\b);}
  \draw[->,thick] (m)--(xl);
  \foreach \a/\b in {xl/cs,cs/g,g/t,t/ms}{\draw[->,thick] (\a)--(\b);}
  \node[pl] (a1) at (0,-4.5){Ausbaustufe: XML-Erstellungsmodus (Kommandozeile, Hardware-Filter)};
  \node[pl] (a2) at (12.4,-4.5){Ausbaustufe: \texttt{<publish>}-Sektion (Lager-Rückschrieb)};
  \node[pl] (a3) at (9.3,-4.5){Ausbaustufe: dritte Factory LaTeX (Blatt-Ordnung, Messreihe je Binary)};
  \draw[->,dashed] (a1)--(ms);
  \draw[->,dashed] (a2)--(xl);
  \draw[->,dashed] (a3)--(cs);
\end{tikzpicture}
\caption{Die automatisierte Evaluations-Kette von der Anwender-XML bis zum Manuskript. Obere Reihe: Die
XML deklariert Was, Wo und Wann; der Validator prüft sie gegen die Registries; der Planer emittiert den
deterministischen Plan; die CEB baut je System-Permutation die Tier-Binaries; die Messung erhebt die
drei Mess-Ebenen und die Zähler mit einem Mess-Thread. Untere Reihe: Die Mappe ist der Stamm (ein Datenblatt
unter konstantem Blattschlüssel; die Auffächerung in ein Blatt je Unter-Achsen-Permutation ist Soll), die CSV
ihr Kind, die Generatoren rendern je Sprache in das \texttt{tabellen}-Verzeichnis, das Manuskript bindet die
Dateien ein --- jede Zahl trägt ihren Lauf. Gestrichelt: Ausbaustufen (XML-Erstellungsmodus,
\texttt{<publish>}-Sektion, dritte Factory).}\label{fig:eval-xml-kette}
\end{figure}

\subsection{Workloads und Datensätze}\label{sec:eval-workloads}
%% OWNER-VORLAGE: OV-02 --- Default-Workload: die Voll-Permutation über alle verfügbaren Workloads ist
%% Plan-Ist-Stand (dynamische Dimension workload); der Treiber fällt für profil-lose Module auf genau einen
%% Workload zurück. Durchführungs-Stand an dieser einen Stelle; Zweiteilung in beide Richtungen.
Phase~5 der Mess-Pipeline (Execute, Abschnitt~\ref{sec:impl-pipeline}) leitet den YCSB-Workload je
Permutation als Fallback-Heuristik aus dem Traversal-Tag des Profils ab
(Tabelle~\ref{tab:workload-routing}), sofern kein \texttt{<expected\_workload>}-Tag des Profils die
Heuristik überschreibt (Abschnitt~\ref{sec:sota-instances}); das Routing matcht dabei per Teilstring
(\texttt{RANGE}, \texttt{HOT\_PATH}) und nicht über eine geschlossene Tag-Liste, und der Wert des
\texttt{<expected\_workload>}-Tags wird gegen die sechs Token \texttt{YCSB\_A} bis \texttt{YCSB\_F}
aufgelöst. Der globale Default-Workload ist die Permutation über alle verfügbaren Workloads; die
Laufzeit-Dimension \texttt{workload} ist dafür als dynamische Dimension der Mess-Registry
registriert, gespeist aus dem Last-Rahmenwerk YCSB~\cite{cooper2010ycsb}. Durchführungs-Stand: Der
Treiber fällt für profil-lose Module auf genau einen Workload zurück, die Workload-Option des Laufs;
die Voll-Permutation ist Plan-Ist-Stand und wird mit dem Kampagnen-Fenster gebaut. Als
Datengrundlage dienen reale String-Datensätze unterschiedlicher Charakteristik
(Tabelle~\ref{tab:datasets}), jeder mit einer Reproduzierbarkeits-Akte aus Quelle, Prüfsumme,
Zeilenzahl, Vorverarbeitung und Seed-Regel für die synthetischen Miss-, Update- und Delete-Schlüssel;
vier der sechs Datensätze (\texttt{url}, \texttt{protein}, \texttt{tpcds-id}, \texttt{trec-terms})
tragen eine gegen die vorverarbeiteten Dateien berechnete reale Akte, \texttt{dna} die Akte der
Pizza\&Chili-Quelle, und \texttt{xml} trägt eine Spezifikation ohne Prüfsumme, weil die Datei nicht
lokal vorliegt --- eine erfundene Prüfsumme gibt es nicht (n/a). Die Akten liegen als XML-Dateien je
Datensatz im Werkzeug-Repository, und ein Kanon-Test prüft die realen Akten gegen den Sechser-Kanon
der Tabelle.

\begin{table}[!htbp]
  \centering
  \caption{Profil-bewusstes Workload-Routing (Fallback-Heuristik)}\label{tab:workload-routing}
  \begin{tabular}{lll}
    \toprule
    \textbf{Traversal-Tag} & \textbf{YCSB-Workload} & \textbf{Charakteristik} \\
    \midrule
    \texttt{RANGE\_SCAN} / \texttt{RANGE\_HOT\_PATH} & YCSB-E & Bereichsabfragen \\
    \texttt{HOT\_PATH} / \texttt{PRTART\_HOT\_PATH} & YCSB-A & 50/50 Lesen/Aktualisieren, Zipf \\
    \texttt{STANDARD} (Default) & YCSB-C & nur Lesen \\
    \bottomrule
  \end{tabular}
\end{table}

\begin{table}[!htbp]
  \centering\small
  \caption{Reale Datensätze, ihre Charakteristik und ihre Reproduzierbarkeits-Akte}\label{tab:datasets}
  \begin{tabular}{lll}
    \toprule
    \textbf{Datensatz} & \textbf{Charakteristik} & \textbf{Akte} \\
    \midrule
    \texttt{url} & sehr lange Schlüssel, lange gemeinsame Präfixe, mittleres Alphabet & real \\
    \texttt{dna} & kleines Alphabet, hohe Wiederholung, tiefe Präfix-Entropie & Pizza\&Chili-Akte \\
    \texttt{protein} & biologische Sequenzen, kleines bis mittleres Alphabet & real \\
    \texttt{xml} & strukturierte Strings mit wiederkehrenden Präfixen & n/a (Datei nicht lokal) \\
    \texttt{tpcds-id} & datenbanknahe IDs, sortiert und unsortiert & real \\
    \texttt{trec-terms} & Wörterbuch-/Suchindex-Terme mit definierten Miss-Schlüsseln & real \\
    \bottomrule
  \end{tabular}
\end{table}

\subsection{Plattformen, Hersteller-Lanes und Regimes}\label{sec:eval-platforms}
%% OWNER-VORLAGE: OV-K5-01 / OV-K3-03 --- Barnard-CPU-Familie mit ZIH-Beleg (Intel Sapphire Rapids, Xeon
%% Platinum 8470) gegen den Code-Kommentar (AMD EPYC 7763); ZIH deklariert, nicht vermessen (n/a),
%% gleichlautend mit Kapitel 3. OV-02 --- XML-Erstellungsmodus als Plan-Ist-Stand (Verweis Kapitel 4).
%% Defaults gesetzt; Widerspruch genügt.
Die Messungen zielen auf die registrierte Flotte des Apparats: prod1, die AMD-Maschine mit einem
Ryzen~9~9950X3D (Zen~5; Maschinen-Signatur \texttt{prod1\_zen5} mit 22 Merkmalen inklusive AVX-512), % chktex 29
zugleich die Entwicklungs-Workstation der Verifikation; prod2, die Intel-Maschine mit einem Core
i9-12900K (Alder Lake; Signatur \texttt{prod2\_alder\_lake} mit neun Merkmalen, ohne AVX-512, mit P-
und E-Kernen); und ein ODROID-Board mit Intel-Gracemont-Kernen (Signatur \texttt{odroid\_gracemont},
neun Merkmale). Der ZIH-Cluster der TU Dresden ist als Unter-Achse IS3 der Ziel-ISA deklariert
(Barnard, Capella, GraceHopper, N/A); der Barnard-Knoten trägt Intel-Sapphire-Rapids-CPUs, zwei Xeon
Platinum 8470 je Knoten~\cite{zih_hpc_compendium_hardware}, und Capella ist der GPU-Cluster mit
AMD-Genoa-CPUs und NVIDIA H100. In keiner Maschinen-Signatur ist ein ZIH-Knoten registriert, und ein
Lauf dort ist nicht vermessen (n/a) --- gleichlautend mit Abschnitt~\ref{sec:measurement-system}. Eine
HBM-Plattform ist in keiner Signatur und in keiner Achse als Plattform registriert; die Plattform-Sonde
trägt ein HBM-Merkmal und die Allokator-Achse kennt PIM-Sonderfälle, mehr nicht --- sie bleibt
geplant, nicht registriert. Eine \texttt{IPlatformProbe} (Phase~1) meldet je Plattform die lauffähige
Teilmenge der Achsen (etwa AVX-512 nur dort, wo verfügbar), wobei der Ausschluss einer nicht
verfügbaren Erweiterung stets \emph{maschinenseitig} fällt --- mit Warnung und benannter
Fehlerklasse, nie als stiller Profil-Filter, denn das Experiment-Profil bleibt für alle Maschinen
gleich; je Plattform werden ein Scalar- und ein AVX2-Build vorbereitet, ein AVX-512-Build nur dort, wo
die Signatur das Merkmal trägt. Die verfügbaren Achsen-Optionen sind je Plattform über die
Planer-Sonde prüfbar; ein kommandozeilengeführter XML-Erstellungsmodus mit Plattform-Filter ist
Plan-Ist-Stand und wird im Kampagnen-Fenster gebaut (Abschnitt~\ref{sec:impl-pipeline}). Auf
Hybrid-CPUs werden P- und E-Kerne getrennt vermessen statt zusammenaggregiert, auf Intel wie auf AMD\@;
die Messmechanik dafür sind die Core-PMU-Domänen des Kernels (auf Intel die sysfs-Namen
\texttt{cpu\_core} und \texttt{cpu\_atom}), während die Benennung der Unter-Achse vendorneutral bleibt
--- die sysfs-Namen stehen ausdrücklich auf der Verbotsliste der ISA-Unter-Achsen-Etiketten. Die
Kern-Modi folgen dem Sechs-Modus-Modell aus Abschnitt~\ref{sec:axis-triad}: vier Pinning-Modi je
Kern-Sorte (P- oder E-Kerne, einzelner Kern oder alle Kerne der Sorte), der ungepinnte Referenzlauf
und \emph{HybridAware} als sechster Modus (Scheduling-Dimension \texttt{hetero\_core\_dispatch}), der
für die Cache-Line-Awareness-Messung verpflichtend ist; die Verdrahtung des Kern-Pinnings und der
Kern-Modi in das Lauf-Profil ist als Bau-Schritt benannt (Abschnitt~\ref{sec:impl-system-axes}).

Der produktive Mess-Betrieb ist ein \emph{Zwei-Maschinen-Betrieb} --- die AMD-Maschine prod1 neben
der Intel-Maschine prod2 ---, und die zählerbasierte Erhebung führt dafür zwei je einzeln
scharfschaltbare Hersteller-Lanes, \texttt{pmc:amd} auf prod1 und \texttt{pmc:intel} auf prod2, mit je
eigener PMU-Ressourcengruppe; die Lanes laufen absichtlich parallel, weil verschiedene Hosts
verschiedene PMUs tragen und eine gemeinsame Gruppe sie grundlos serialisieren würde. Die
Vollständigkeit ist eine Auswertungs-Regel, kein CI-Gate: Eine Messreihe gilt erst als vollständig,
wenn beide Lanes ihre Zähler-Voraussetzungen nachgewiesen und geliefert haben; eine einseitige
Lane-Fahrt wird als Teil-Messung ausgewiesen, nie still als Ganzes geführt
(Abschnitt~\ref{sec:eval-validity}). Auf den Produktions-Zielmaschinen soll zudem jede Messung unter
\emph{zwei Betriebssystem-Regimes} erhoben werden, um Produktions-Treue von Mess-Tiefe zu trennen: Ein
\emph{immutables} Betriebssystem (Talos) spiegelt den Produktiv-Einsatz und erzwingt
boot-/übersetzungsstatische Cache-Einstellungen (Abschnitt~\ref{sec:hardware-history}); ein
\emph{root-Linux mit vollem Hardware-Zähler-Zugriff} (\texttt{perf} und die modellspezifischen
Register, MSR) eröffnet hingegen überhaupt erst den Zugang zu den zählerbasierten Mess-Kategorien
(Abschnitt~\ref{sec:sota-instances}) und zu den laufzeit-konfigurierbaren cache-nahen
Achsen-Parametern (Prefetch-Distanz, Hardware-Prefetcher). Das zweite Regime ist \emph{entworfen,
nicht vollzogen}: Die berichteten Messungen stammen sämtlich vom root-Linux-Regime der beiden
Produktionsmaschinen; der Talos-Spiegel bleibt als Anforderung bestehen.

\subsection{Mess-Ebenen, Kollektoren und Zähler}\label{sec:eval-counters}
Davon zu trennen ist, welche Zähler dieser Zugang tatsächlich füllt
(Tabelle~\ref{tab:eval-pmc-zaehler}): Die Mess-Kette öffnet über
\texttt{perf\_event\_open}~\cite{linux_perf_event_open,demelo2010perf} vier generische Zähler ---
L1D-Misses, Last-Level-Misses, dTLB-Misses und Branch-Misses
(\texttt{PERF\_COUNT\_HW\_BRANCH\_MISSES}, portabel auf Intel und AMD). Für L2-Misses und
Kohärenz-Invalidierungen gibt es keinen portablen generischen Zähler; beide werden modell-gebunden über
einen RAW-Event-Katalog erhoben, dessen Kodierungen am Objekt kreuzgeprobt sind: Für
Zen~5\footnote{Zen~5: die AMD-Mikroarchitektur des Ryzen~9~9950X3D (Granite Ridge, CPUID-Familie 26); % chktex 29
prod1 trägt zwei ungleich große L3-Komplexe.} ist \texttt{l2\_cache\_req\_stat.ic\_dc\_miss\_in\_l2} das
Event \texttt{0x64} mit Maske \texttt{0x09} (Demand-Misses in L2 ohne L2-Prefetch) und % chktex 29
\texttt{ls\_dmnd\_fills\_from\_sys.remote\_cache} das Event \texttt{0x43} mit Maske \texttt{0x14} % chktex 29
(Demand-Fills aus dem Cache eines anderen CCX, eines Kern-Komplexes mit eigenem L3); ein wörtlicher
Invalidierungs-Zähler existiert im Core-Event-Satz nicht, die Spalte trägt genau diese Belegung. Die
Kodierungen stammen aus der Eventtabelle des Linux-Kernels für Zen~5~\cite{linux_perf_amdzen5}, aus
der \texttt{perf} die benannten Events bindet. Modelle ohne Katalog-Eintrag tragen n/a, nie eine
geratene Kodierung; für Intel liegt kein Eintrag vor, weil die Kreuzprobe auf einem Intel-Host
aussteht. Das Last-Level-Event scheitert auf Zen~5 beim Öffnen (\texttt{ENOENT}). IPC/CPI ist als
Kategorie deklariert, aber nicht verdrahtet --- dem Zähler-Verbund fehlen die Quellen für Instruktionen
und Zyklen ---, und die Energie-Aufnahme wird best-effort aus der RAPL-Sysfs-Zone\footnote{RAPL
(Running Average Power Limit): die Energie-Zähler der Prozessoren, die das Betriebssystem über die
Powercap-Schnittstelle in Mikrojoule ausliest~\cite{david2010rapl,khan2018rapl}; ohne Leserecht auf
die Zone bleibt der Wert leer.} gelesen und bleibt ohne Zonen-Leserecht leer. Die Ehrlichkeits-Regel
ist als Feld-Flag realisiert: Je Zähler trägt der Zähler-Verbund ein
\texttt{*\_source\_available}-Flag, sieben Flags für sieben Zähler; eine Zahl bedeutet, dass die
Quelle geliefert hat --- auch eine echte Null ---, n/a bedeutet, dass die Quelle fehlt, und diese
Unterscheidung reist Ende-zu-Ende bis in Mappe, Tabellen und Diagramme. Hardware-Zähler gelten dabei
nur nach Kreuzprobe als vertrauenswürdig~\cite{weaver2008counters}; die Kreuzprobe für Zen~5 ist am
Objekt geführt, die für Intel steht aus.
\begin{table}[htbp]\centering\footnotesize
\begin{tabular}{@{}l>{\raggedright\arraybackslash}p{4.6cm}>{\raggedright\arraybackslash}p{3.2cm}>{\raggedright\arraybackslash}p{2.3cm}@{}}
\toprule
Kategorie (Spalte) & Quelle & Intel-Stand & Ehrlichkeits-Flag \\
\midrule
\texttt{cache\_misses\_l1} & generisches perf-Event (L1D-Read-Miss) & generisch: ja & eigenes Flag \\
\texttt{cache\_misses\_l2} & RAW-Katalog Zen~5: Event \texttt{0x64}, Maske \texttt{0x09} (Demand-Misses in L2) & kein Katalog-Eintrag, Kreuzprobe ausstehend (n/a) & eigenes Flag \\ % chktex 29
\texttt{cache\_misses\_l3} & generisches perf-Event (Last-Level); scheitert auf Zen~5 beim Öffnen (n/a) & generisch: ja & eigenes Flag \\
\texttt{dtlb\_misses} & generisches perf-Event (dTLB-Read-Miss) & generisch: ja & eigenes Flag \\
\texttt{branch\_misses} & generisches perf-Event (\texttt{PERF\_COUNT\_\allowbreak{}HW\_\allowbreak{}BRANCH\_\allowbreak{}MISSES}) & generisch: ja & eigenes Flag \\
\texttt{coherence\_invalidations} & RAW-Katalog Zen~5: Event \texttt{0x43}, Maske \texttt{0x14} (Fremd-CCX-Fills) & kein Katalog-Eintrag, Kreuzprobe ausstehend (n/a) & eigenes Flag \\ % chktex 29
\texttt{energy\_micro\_joules} & RAPL-Sysfs-Zone, best-effort (Leserecht) & wie AMD (Leserecht) & eigenes Flag \\
IPC/CPI & deklariert, nicht verdrahtet (Quelle fehlt) & n/a & --- \\
\bottomrule
\end{tabular}
\caption{Die Zähler der Mess-Kette und ihre Quellen: vier generische perf-Events (portabel), zwei
modell-gebundene RAW-Events aus dem Katalog für Zen~5 (kreuzgeprobt am Objekt), die RAPL-Energie als
best-effort-Lesung und IPC/CPI als deklarierte, nicht verdrahtete Kategorie. Je Zähler trägt der
Zähler-Verbund ein \texttt{*\_source\_available}-Flag; eine Zahl heißt \enquote{Quelle hat geliefert},
n/a heißt \enquote{Quelle fehlt}. Der Intel-Stand nennt, was ohne Katalog-Eintrag auf prod2
erhoben wird.}\label{tab:eval-pmc-zaehler}
\end{table}

Zeit- und Beobachter-Messung sind im Mess-Checkpoint (\texttt{checkpoint\_measure}) vereint; ihre
Verantwortlichkeiten trennt die Mess-Tooling-Haupt-Achse in genau dieser Reihenfolge: \emph{wallclock}
(w) --- die Wall-Clock der Gesamtläufe im CacheEngineBuilder über die Last-Rahmenwerke,
parameter-leer; \emph{macro} (ma) --- die Zeit der Genus-Interface-Funktionen der Tier-Binary;
\emph{micro} (mi) --- die Zeit und die Erfolgs-Parameter je Achsen-Interface
(Abschnitt~\ref{sec:measurement-system}, Abbildung~\ref{fig:traeger-flaechen}); Teilmengen der drei
Ebenen sind erlaubt, jede andere Reihenfolge ist syntaktisch falsch, und w/ma/mi ist der registrierte
Alias. Drei Kollektoren tragen die Erhebung als Bausteine der Mess-Achse \texttt{collector} ---
Wall-Clock, Beobachter-Schnappschuss und PMC ---, und die sechzehn Mess-Kategorien sind die
Unter-Achse der Tooling-Achse, nicht ihre Haupt-Auffächerung. Die Auswertung hält die drei Ebenen
getrennt: Achsen-Interface, Gattungs- und Genus-Interface und die Wall-Clock des CacheEngineBuilders
werden nie gemischt (Abschnitt~\ref{sec:eval-sensitivity}); PMC-Messungen als eigener Abbildungstyp des
Anhang-Generators sind eine Ausbaustufe. Durchführungs-Stand: Registry und Ebenen-Kanon sind gebaut;
die Ordnungs-Validierung an den drei Serialisierungs-Stellen und der Kollektor als eigenständiges
Bauteil sind als Posten benannt.

Die zeit- und beobachter-basierten Kategorien laufen unter beiden Regimes identisch und sichern die
Vergleichbarkeit; die PMC-Kategorien werden nur unter dem privilegierten Regime erhoben. Je Kategorie
trägt die Registry das Regime als Ordinal (0 für Zeit und Beobachter, 1 für Zähler): neun
Zeit-/Beobachter-Kategorien und sieben Zähler-Kategorien. Betriebssystem, Kernel und Zähler-Domäne sind
je Plattform zu protokollieren: Die Betriebssystem-Sonde erhebt Kernel und Kennung des Systems als
System-Achse, und die Zähler-Quelle hält je Zähler fest, auf welcher Core-PMU-Domäne er geöffnet hat
(generisch oder, auf Hybrid-Maschinen, der sysfs-Name der Domäne); eine eigene Spalte der Mess-Ausgabe
tragen diese drei Angaben im WIDE-Schema nicht --- die Spalten \texttt{platform} und
\texttt{build\_version} tragen Plattform-Tag und Bau-Version je Zeile, die Protokoll-Spalte je
Domäne ist Anforderung.

\subsection{Bau-Umfang und Optimierungsstufen}\label{sec:eval-build-scope}
%% OWNER-VORLAGE: OV-02 --- O0 bis O3 als Kunden-Wahl in der XML samt Durchreichung weiterer Compile-Flags:
%% Verfahren als Plan-Ist-Stand, Durchführungs-Stand an dieser einen Stelle (Registry führt fünf Stufen,
%% das Profil deklariert O2 und O3). Default gesetzt; Widerspruch genügt.
Der Umfang des vollständigen Baus ist als \emph{Soll} festgelegt. Das Profil deklariert die
Bau-Steuergrößen: zwei Optimierungsstufen (O2, O3) mal zwei SIMD-Stufen (keine Erweiterung, AVX2), also
vier Bau-Varianten je Permutation (Abschnitt~\ref{sec:axis-triad}). Beide Stufen sind in der
System-Registry als IEEE-754-deterministisch deklariert; die Registry führt fünf Stufen (O0 bis O3 und
Ofast), und ausgeschlossen aus dem Profil ist allein Ofast, weil es den Gleitkomma-Determinismus nach
IEEE~754 und den Run-to-Run-Determinismus des Apparats sowie den CRC64-Anker des golden-Katalogs
bräche --- die Registry markiert genau diese Stufe als nicht deterministisch. Die maximale
Optimierung bleibt damit wählbar: Der Haus-Standard ist global O2, O3 wird auf XML-Eingangswunsch
mitgebaut, und die gewählte Stufe ist über das Werkzeugketten-Glied des Fingerprints
identitätswirksam, sodass O2- und O3-Bestände nie stillschweigend als Cache-Treffer verwechselt
werden; O0 bis O3 als Kunden-Wahl in der XML samt Durchreichung weiterer Compile-Flags sind
Plan-Ist-Stand und werden im golden-Fenster gebaut. Je Bau-Variante wird der golden-Referenzkatalog
gebaut: $2^{17} = 131072$ Binary-Identitäten, weil siebzehn Organ-Haupt-Achsen je zwei Vertreter
freischalten und die achtzehnte (\texttt{persistence\_target}) gepinnt ist
(Abschnitt~\ref{sec:mess-chain}); die System-Achsen multiplizieren die Bau-Matrix, nie die
\texttt{binary\_id}. Die Mächtigkeit des realen Permutationsraums rechnet der Planer dynamisch aus
seinen Eingängen (Subkommando \texttt{simulate}); sie ist kein statischer Nenner. Ergebnisse der
Voll-Bau-Läufe liegen nicht vor (Abschnitt~\ref{sec:eval-available}).

\subsection{Fairness und Reproduzierbarkeit}\label{sec:fairness}
%% OWNER-VORLAGE: OV-K5-04 --- Label sec:fairness in Kapitel 5 gesetzt, damit die Code-Verweise des
%% Profil-Validators und der Mess-Achse wieder auf das gebaute Dokument zeigen (Alternative: Code-Nachzug
%% auf sec:eval-method). Default gesetzt; Widerspruch genügt.
Faire Vergleichbarkeit ist über feste Regeln gesichert. Jeder Vergleich trennt einen gemeinsamen
Minimalmodus (Common-Denominator: externe Werte-Handles, keine PRT-ART-Spezialpfade) vom
PRT-ART-Native-Modus (Inline-Umschaltung und adaptive Seitentyp-Wahl je Verzweigungspunkt); die
Anwender-XML deklariert den Modus je SOTA-Reihe als \texttt{fairness\_mode} mit den zwei erlaubten
Werten \texttt{common\_denominator} und \texttt{native}, und der Validator weist jeden anderen Wert ab.
Durchführungs-Stand: Der Modus ist Reihen-Metadatum ohne Mess-Einfluss --- er reist als Tag-Spalte der
Mess-Zeile und im Wiederaufnahme-Stempel, nie in der \texttt{binary\_id} ---, die Kompositions-Pinnung
des Minimalmodus ist daten-gebunden, und zwei Reihen desselben Prüflings mit verschiedenem Modus
erzeugen deshalb dieselbe \texttt{binary\_id}, was der Validator als Warnung meldet. Alle Verfahren
erhalten identische, normalisierte Byte-Schlüssel sowie identische Miss-, Update- und Delete-Mengen;
der Warmup ist der Zwei-Phasen-Treiber des Lastprofil-Orchestrators (sichern, aufwärmen,
zurückrollen, messen), sonst wird kalt gemessen, und jede Konfiguration wird in mehreren getrennten
Läufen vermessen (Wiederholungs-Dimension \texttt{repetition}). Thread-Pinning: Der Aktuator --- ein
bereichsgebundenes Pinning des Mess-Threads --- ist gebaut; die Kern-Wahl innerhalb des
Kern-Komplexes ist Aktuator-Politik ohne Emissions-Vertrag, und die Mess-Emission des Planers
deklariert das Pinning als Soll und schreibt den Ist-Zustand \emph{ungepinnt} daneben, weil der
Aktuator im Lauf-Profil nicht verdrahtet ist --- auf prod1 mit seinen zwei ungleich großen
L3-Komplexen (96 und 32~MiB) ist ein ungepinnter Lauf nicht reproduzierbar, was die Emission wörtlich
sagt. Fehlende Fähigkeiten einer Baseline --- etwa Bereichs-Scans einer Hash-Tabelle --- werden als
n/a ausgewiesen statt verdeckt emuliert; die Mechanik dafür sind die Stichproben-Status
\emph{NotApplicable}, \emph{SourceUnavailable} und \emph{Failed} der Mess-Achsen, sodass fehlende
Fähigkeiten und fehlende Quellen n/a tragen, nie eine Null. Provenienz: Werkzeugkette
(Compiler-Kennung und -Version), Flags, ISA-Zusammenfassung und die Commit-Stände der vier
Repositories werden zur Konfigurationszeit als Bau-Provenienz erfasst und in das Provenienz-Manifest
der Host-Exporte geschrieben; der Allokator ist als Organ-Achse Teil jeder \texttt{binary\_id}; eine
Provenienz je einzelner Fremdbibliothek --- ihr eigener Compiler, ihre Flags, ihr Commit-Hash --- ist
zu protokollieren (Anforderung). Die Latenz-Perzentile bestimmt der Auswertungs-Baum der Cache-Engine
nearest-rank aus den Rohwerten --- Hyndman und Fan, Typ~1: Rang $k = \lceil q \cdot n \rceil - 1$,
ohne Interpolation~\cite{hyndman1996quantiles} ---, sodass jedes Perzentil ein real gemessener Wert
ist\footnote{Nearest-Rank: Das $q$-Quantil ist der kleinste Stichprobenwert, unter dem mindestens der
Anteil $q$ der Stichprobe liegt; interpolierende Verfahren (Typ~7) erzeugen dagegen Werte, die nie
gemessen wurden.}; die Darstellungs-Werkzeuge des Werkzeug-Repositories tragen eigene Kopien der
Umrechnung (eine benannte Lücke), und das HDR-Histogramm~\cite{tene_hdrhistogram} ist ein eigenes
Verfahren der Verteilungssicht (Abschnitt~\ref{sec:impl-system-axes}), keine Ergänzung des
Nearest-Rank. Layout- und Umgebungs-Effekten, die ohne erkennbaren Fehler falsche Schlüsse
erzeugen~\cite{mytkowicz2009wrongdata}, begegnet der Apparat mit identischen Schlüsseln, identischer
Bau-Matrix je Variante und der Wiederholungs-Dimension; eine randomisierte Layout-Stabilisierung im
Sinne von Curtsinger und Berger~\cite{curtsinger2013stabilizer} ist nicht Teil des Apparats und wird
in Abschnitt~\ref{sec:eval-validity} als Grenze geführt.

%% OWNER-VORLAGE: OV-02 --- Original-Compiler-Doktrin: Nachweis-Seite (is_original-Markierung, Validator,
%% Relock-Gate) gebaut; automatischer Einzug der vendorten Quellen mit ihrem Original-Compiler als
%% Plan-Ist-Stand. Durchführungs-Stand an dieser einen Stelle; Widerspruch genügt.
Für fremde Bibliotheken gilt die Original-Compiler-Doktrin: Sezierte Achsen-Funktionen werden zur
Übersetzungszeit eingebunden und zu einem wieder originalen Binary-Block kompiliert, sodass jeder
sezierte Teilalgorithmus seine Leistungseigenschaften behält; wo umsetzbar oder vergleichbar, wird die
Fremdquelle mit ihrem Original-Compiler gebaut. Der Nachweis der Originalität ist gebaut: die
\texttt{is\_original}-Markierung je Funktion (ein Makro der Cache-Engine samt Mixin-Basis der
Achsen-Bausteine), der Validator als eigenes Programm und das manuelle Relock-Gate der CI, das die
Lock-Dateien der Markierungen neu verankert; der Beleg je Baustein steht in Anhang~\ref{app:blocks}.
Die H2-Akte markiert Rekonstruktionen ohne vendorte Originalquelle als n/a (22 von 33,
Abschnitt~\ref{sec:eval-hypotheses}). Durchführungs-Stand: Der automatische Einzug der vendorten Quellen
in den Bau mit ihrem Original-Compiler ist Plan-Ist-Stand und liegt im Fenster nach dem Trigger.

%% OWNER-VORLAGE: OV-02 --- Lager im Mess-Betrieb: Bestand 1/2 und Schlüssel-Schema gebaut; Bestand 3/4,
%% Planer-Blattfunktion, CEB-Vollplatzierung, Hybrid-Ablage, Commit-Skip, rebuild-Flag und XML-<publish>
%% als Lager-Vollausbau benannt. Durchführungs-Stand an dieser einen Stelle; Widerspruch genügt.
Auch der Mess-\emph{Betrieb} ist Teil dieser Fairness-Regeln, und er läuft über das Lager. Das Lager
ist der Bestand aller gebauten Tier-Binaries und ihrer Messwerte, adressiert über
Übersetzungszeit-Schlüssel (Abschnitt~\ref{sec:lager}; Abbildung~\ref{fig:lager-bestaende} zeigt seine
Struktur): Bestand~1 trägt die Binary unter ihrem SHA-512-Fingerprint, Bestand~2 die Messwerte unter
dem Paar aus dem Fingerprint der vermessenen Binary und der Hardware-Identität der messenden Maschine;
die Zelle aus ISA, Optimierung und Erweiterung steht daneben geklammert, nie im Digest. Weil der
Mess-Gate-Zustand ein Fingerprint-Glied ist, liegen die Binary mit und die Binary ohne Messfühler als
zwei Einträge des Bestands~1 nebeneinander. Die Messung ist immer eingeschaltet; im Lager wird nur
nachgemessen, was dort fehlt: Eine vorgefundene Binary wird übersprungen, wenn ihr Fingerprint bekannt
und ihr grünes Prüf-Testat am Binary inventarisiert ist, sonst wird gebaut oder gemessen
(Abbildung~\ref{fig:eval-lager-kreislauf}). Die CEB steuert die Parallelität, indem sie parallele
Anfragen an die Tier-Binaries erzeugt --- die Ausführungs-Threads entstehen in der CEB, in einem
Worker-Pool mit deterministischer Ergebnis-Reihenfolge ---, während das Threading der Gattung in der
Tier-Binary bleibt. Durchführungs-Stand: Bestand~1 und~2 sowie das Schlüssel-Schema aller vier
Bestände sind gebaut; Bestand~3 und~4 (synthetisierte Kurven als Dateien, Lösungen je Anwender-XML),
die Planer-Blattfunktion, die Vollplatzierung der CEB, die Hybrid-Ablage, der Commit-Skip, das
rebuild-Flag und die XML-Sektion \texttt{<publish>} sind der Lager-Vollausbau.
\begin{figure}[htbp]\centering
% Kein \resizebox: vier Kacheln im Raster 4.1cm (Gesamtbreite ~15.3cm) passen nativ in \textwidth;
% Regel: resizebox nur zum VERKLEINERN (vgl. fig:m-model).
\begin{tikzpicture}[font=\footnotesize,>=Stealth,thick,
  k/.style={draw,rounded corners,fill=blue!7,text width=2.7cm,minimum height=1.15cm,align=center,inner sep=2pt,font=\scriptsize},
  lg/.style={k,fill=green!10},
  dc/.style={draw,fill=orange!14,text width=2.7cm,minimum height=1.15cm,align=center,inner sep=2pt,font=\scriptsize},
  pl/.style={k,fill=orange!10,dashed,font=\tiny,minimum height=0.9cm}]
  \node[k]  (bau)  at (0,0){\textbf{Bau (CEB)}\\Batches zu je 4096 Binaries, inventarisierte Wiederaufnahme};
  \node[lg] (b1)   at (4.1,0){\textbf{Bestand 1: Binary}\\Schlüssel: SHA-512-Fingerprint; Zelle daneben geklammert};
  \node[dc] (skip) at (8.2,0){\textbf{Skip-Kriterium}\\Fingerprint bekannt \emph{und} Prüf-Testat grün inventarisiert?};
  \node[k]  (mess) at (12.3,0){\textbf{Messung}\\Prüf-Dock; ein Mess-Thread; die CEB erzeugt parallele Anfragen};
  \node[lg] (b2)   at (12.3,-2.6){\textbf{Bestand 2: Messwerte}\\Schlüssel: Fingerprint $+$ Hardware-Identität; Zelle daneben};
  \node[lg] (mp)   at (8.2,-2.6){\textbf{Mappe (xlsx)}\\Projektion csv; Anhang-Generatoren};
  \node[k]  (ref)  at (4.1,-2.6){\textbf{Struktur des Lagers}\\Abbildung~\ref{fig:lager-bestaende}: Baum, Bestände 1--4, Schlüssel-Schema};
  \node[pl] (vp)   at (0,-2.6){Lager-Vollausbau: Bestand 3/4, Hybrid-Ablage, Commit-Skip, rebuild-Flag, \texttt{<publish>}};
  \draw[->] (bau)--(b1);
  \draw[->] (b1)--(skip);
  \draw[->] (skip)--node[above,font=\tiny]{nein: messen}(mess);
  \draw[->] (mess)--(b2);
  \draw[->] (b2)--(mp);
  \draw[->] (skip)--node[right,font=\tiny,align=left]{ja: überspringen,\\Messwerte liegen}(mp);
  \draw[->] (skip.north) -- ++(0,0.75) -| node[pos=0.25,above,font=\tiny]{Fingerprint unbekannt: bauen} (bau.north);
  \draw[->,dashed] (vp)--(ref);
\end{tikzpicture}
\caption{Der Lager-Kreislauf im Mess-Betrieb: Die CEB baut in Batches zu je 4096 Binaries; jede Binary
liegt unter ihrem SHA-512-Fingerprint in Bestand~1; das Skip-Kriterium überspringt, was mit grünem
Prüf-Testat inventarisiert ist, und lässt sonst bauen oder messen; die Messung läuft am Prüf-Dock mit
einem Mess-Thread, während die CEB die parallelen Anfragen erzeugt; die Messwerte liegen unter
Fingerprint und Hardware-Identität in Bestand~2 und fließen in die Mappe. Die Struktur des Lagers
(Baum, vier Bestände, Schlüssel-Schema) zeigt Abbildung~\ref{fig:lager-bestaende}; gestrichelt der
Lager-Vollausbau.}\label{fig:eval-lager-kreislauf}
\end{figure}

Die Messung bleibt strikt vom Bau getrennt: Mess-Läufe starten nie automatisch mit einem grünen Bau,
sondern werden ausdrücklich ausgelöst --- ein grüner Bau beweist Bau-Fähigkeit, kein Mess-Ergebnis.
Die Ablaufmethodik kennt vier Arbeitsmodi --- \texttt{build}, \texttt{measure}, \texttt{compare} und
\texttt{release} (Abschnitt~\ref{sec:impl-pipeline}) ---; der Mess-Modus misst mit genau \emph{einem}
Mess-Thread, eine parallele Mess-Ausführung ist über keine Registry-Zeile scharfgeschaltet und nur über
die Zustands-Injektion des Laufeintrags erreichbar, und \texttt{--debug} ist ein Schalter der
Planer-Shell für den Verdrahtungs-Check, der für Anwender gesperrt bleibt, kein Arbeitsmodus; das
Threading der Gattung selbst bleibt dabei der Tier-Binary überlassen, statt vom Harness überschrieben
zu werden. Ein im Lager vorgefundenes Binary wird nur dann übersprungen, wenn sein grünes Prüf-Testat
am Binary inventarisiert ist (Abbildung~\ref{fig:eval-lager-kreislauf}), und vor jeder Messung steht
die Konformitätsprüfung am Prüf-Dock (Abschnitt~\ref{sec:impl-contracts}): das Konformitäts-Gate des
Genus-Vertrags vor den Post-Compile-Testaten der Mess-Interfaces, daneben das Konsistenz-Gate der
Messung; Bau- wie Mess-Läufe werden je Maschine in Meilenstein-Batches von je 4096 Binaries am Stück
abgearbeitet --- das Korn ist eine Übersetzungszeit-Konstante, deren Abweichung ein Übersetzungsfehler
ist ---, mit gleichverteilter Zuteilung der gemeinsam baubaren Binaries über die Maschinen und
inventarisierter Wiederaufnahme nach Unterbrechung.

\subsection{Der heuristische Gegenbeweis als vierte Dimension}\label{sec:eval-counterproof}
%% OWNER-VORLAGE: OV-02 / OV-K5-02 --- Hybrid-Stufe: Klassifikation, Dock-Schicht, XML-Parser und Proxy
%% gebaut, Router und Eviction nicht (Durchführungs-Stand an dieser einen Stelle); der Gegenbeweis steht
%% neben der Rangbildung der Sensitivitätsanalyse, nicht in ihr. Defaults gesetzt; Widerspruch genügt.
Über die drei Hypothesen hinaus erhält die Auswertung eine vierte, eigenständige
\emph{Evaluations-Dimension}: den heuristischen Gegenbeweis. Sein Konzept ist in
Kapitel~\ref{ch:measurement} entwickelt (Abschnitte~\ref{sec:heuristics-modes}
und~\ref{sec:heuristics-curves}): Die CacheEngineBuilder-Stufe kompiliert die heuristisch geschätzte
Rekombination der Tier-Binaries zu einer Hybrid-Tier-Binary --- der Hybrid-Modul-Emitter liegt im
Builder-Baum der CEB ---, und die Hybrid-Tier-Binary ist eine eigene Träger-Stufe \emph{hinter} der
CEB, die am Prüf-Dock andockt und selbst keinen Code erzeugt. Das Messkurven-Typsystem liefert die
Schätzgrundlage (Abbildung~\ref{fig:eval-messkurven}): je Achse eine Kurve über ihren Parametern und je
Mess-Ebene wallclock, macro und micro, daraus die Break-Even-Punkte und die geschätzte Rekombination;
die Messfehler-Erkennung formuliert Abschnitt~\ref{sec:heuristics-curves}. Die Bewertungsfrage dieser
Dimension lautet: \enquote{Ist die heuristisch geschätzte Rekombination der Tier-Binaries schneller
als die existierenden Paper-Algorithmen?} Anders als H1--H3, die Aussagen über den Entwurfsraum
prüfen, prüft der Gegenbeweis den Apparat der Arbeit selbst: ob die aus den Messkurven geschätzte
Rekombination im realen Wall-Clock-Verhalten hält, was die Kurven versprechen; er steht deshalb
\emph{neben} der Rangbildung der Sensitivitätsanalyse (Abschnitt~\ref{sec:eval-sensitivity}), nicht
in ihr. Durchführungs-Stand: Klassifikations-Schicht, Dock-Schicht, XML-Parser und Binary-Proxy der
Hybrid-Stufe sind gebaut (sechs der neun geplanten Dateien, Abschnitt~\ref{sec:impl-hybrid-lager});
Router und Eviction brauchen echte Messkurven und sind nicht gebaut.
\begin{figure}[htbp]\centering
% Kein \resizebox: Kurvenfeld 4.8cm + drei Kacheln im Raster 3.0cm (rechte Kante ~13.9cm) und drei
% Varianten-Kacheln zu je 4.6cm (~14.7cm) passen nativ in \textwidth; resizebox nur zum VERKLEINERN.
\begin{tikzpicture}[font=\footnotesize,>=Stealth,
  k/.style={draw,rounded corners,fill=blue!7,text width=2.55cm,minimum height=1.1cm,align=center,inner sep=2pt,font=\scriptsize},
  hy/.style={k,fill=orange!14},
  pl/.style={k,fill=orange!10,dashed,font=\tiny,minimum height=0.8cm},
  bv/.style={draw,rounded corners,fill=green!9,text width=4.5cm,minimum height=1.0cm,align=center,inner sep=2pt,font=\tiny}]
  \draw[->] (0,0) -- (4.8,0) node[below left,font=\scriptsize]{Achsen-Parameter (etwa Working-Set)};
  \draw[->] (0,0) -- (0,3.2) node[above,font=\scriptsize]{Metrik je Ebene};
  \draw[blue!70!black,thick] plot[smooth] coordinates {(0.3,0.5) (1.3,0.8) (2.3,1.5) (3.3,2.3) (4.3,2.9)};
  \draw[green!50!black,thick] plot[smooth] coordinates {(0.3,1.6) (1.3,1.5) (2.3,1.5) (3.3,1.6) (4.3,1.8)};
  \draw[blue!70!black,thin,dashed] plot[smooth] coordinates {(0.3,0.7) (1.3,1.0) (2.3,1.7) (3.3,2.5) (4.3,3.1)};
  \draw[green!50!black,thin,dashed] plot[smooth] coordinates {(0.3,1.8) (1.3,1.7) (2.3,1.7) (3.3,1.8) (4.3,2.0)};
  \fill[red!70!black] (2.3,1.5) circle (2pt);
  \node[font=\tiny,text=red!70!black,anchor=south east] at (2.25,1.6){Break-Even};
  \node[font=\tiny,text=blue!70!black,anchor=west] at (3.35,2.9){Variante A};
  \node[font=\tiny,text=green!50!black,anchor=west] at (3.55,1.35){Variante B};
  \node[font=\tiny,align=left,anchor=north west] at (0.15,3.15){je Achse eine Kurve;\\Bänder wallclock / macro / micro};
  \node[k]  (re) at (6.7,1.6){\textbf{Geschätzte Rekombination}\\Break-Even je Achse und Ebene; Kurven aus dem Lager};
  \node[k]  (em) at (9.7,1.6){\textbf{CEB}\\Hybrid-Modul-Emitter im Builder-Baum: kompiliert};
  \node[hy] (hb) at (12.7,1.6){\textbf{Hybrid-Tier-Binary}\\Stufe hinter der CEB; dockt am Prüf-Dock an; erzeugt keinen Code};
  \node[k]  (wc) at (12.7,-0.6){\textbf{Wall-Clock-Vergleich}\\primär gegen Paper-Rekonstruktionen, sekundär gegen das beste Einzel-Binary};
  \node[pl] (ro) at (9.7,-0.6){Router und Eviction: nicht gebaut (brauchen echte Messkurven)};
  \draw[->,thick] (4.8,1.6)--(re);
  \draw[->,thick] (re)--(em);
  \draw[->,thick] (em)--(hb);
  \draw[->,thick] (hb)--(wc);
  \draw[->,dashed] (ro)--(hb);
  \node[bv] (v1) at (2.35,-2.6){\textbf{Variante 1}: Beobachter in allen Tier-Binaries und im Hybrid-Tier\\Gate-Felder: Tier \texttt{tm}/\texttt{tmi} an, Hybrid \texttt{hm}/\texttt{hmi} an};
  \node[bv] (v2) at (7.35,-2.6){\textbf{Variante 2}: Beobachter nur im Hybrid-Tier\\Gate-Felder: Tier aus, Hybrid an};
  \node[bv] (v3) at (12.35,-2.6){\textbf{Variante 3}: Beobachter in keinem Kompilat, nur Wall-Clock\\Gate-Felder: alle aus};
  \node[font=\tiny,align=center] at (7.35,-3.55){Unsicherheitsband aus den Differenzen 1$\to$2 (Mess-Organe der Tiere) und 2$\to$3 (Mess-Kommandos des Hybrid-Tiers); Gültigkeitskriterium $t_1 \ge t_2 \ge t_3$};
\end{tikzpicture}
\caption{Das Messkurven-Typsystem im heuristischen Gegenbeweis. Links: je Achse eine Kurve der
Metrik über dem Achsen-Parameter, getrennt nach Mess-Ebene (wallclock, macro, micro); der Schnittpunkt
zweier Varianten ist der Break-Even-Punkt, aus dem die Rekombination geschätzt wird. Mitte und rechts:
Die CEB kompiliert die geschätzte Rekombination über den Hybrid-Modul-Emitter; die Hybrid-Tier-Binary
ist die Träger-Stufe hinter der CEB, dockt am Prüf-Dock an und erzeugt selbst keinen Code; verglichen
wird an der Wall-Clock, primär gegen die Paper-Rekonstruktionen, sekundär gegen das beste
Einzel-Binary. Unten: die drei Bau-Varianten des Messfehler-Prüfplans, auseinandergehalten über die
Gate-Felder des Fingerprints. Gestrichelt: Router und Eviction (nicht
gebaut).}\label{fig:eval-messkurven}
\end{figure}

Der Prüfplan hierfür ist zweistufig festgelegt; die zugehörigen Messläufe stehen aus
(Abschnitt~\ref{sec:eval-available}), der Plan gilt im Präsens. Die erste Stufe quantifiziert den
Messfehler, den die Mess-Instanzen selbst eintragen, denn die Schätzgrundlage der Heuristik sind
\emph{beobachtete} Kurven, ihr Anspruch aber gilt dem unbeobachteten Produktions-Verhalten. Dazu wird
dasselbe Experiment in drei Bau-Varianten aufgelegt:
\begin{enumerate}
  \item Beobachter in \emph{allen} Tier-Binaries \emph{und} im Heuristik-Hybrid-Tier,
  \item Beobachter \emph{nur} im Heuristik-Tier-Binary,
  \item Beobachter in \emph{keinem} Kompilat --- gemessen wird nur die Wall-Clock.
\end{enumerate}
Auseinandergehalten werden die Varianten über das Mess-Gate-Glied des Fingerprints
(Abschnitt~\ref{sec:stamp-model}): Es trägt eigene Felder für die Tier-Ebene (\texttt{tm},
\texttt{tmi}) und die Hybrid-Ebene (\texttt{hm}, \texttt{hmi}); der Aus-Zustand aller Felder lautet
\texttt{mg=m0;s0;st0;x0;tw0;tm0;tmi0;hm0;hmi0}, und weil das Glied in den Digest eingeht, liegen die
drei Varianten als drei Binaries mit drei Fingerprints im Lager nebeneinander. Die Differenzen zwischen
den drei Varianten beziffern den Beobachter-Overhead getrennt nach seinem Ursprung --- Mess-Organe der
untergebenen Tiere (Differenz von Variante~1 zu~2) gegen Mess-Kommandos des Hybrid-Tiers (Differenz
von Variante~2 zu~3) --- und ergeben je Workload und Plattform ein \emph{Unsicherheitsband},
unterhalb dessen zwei Konfigurationen als nicht unterscheidbar gelten; das Band wird aus
Wiederholungen mit Konfidenzintervall gebildet, nicht aus einem Einzellauf~\cite{georges2007rigorous}.
Gültigkeitskriterium der Stufe: Mit Messfühler muss langsamer sein als ohne --- für die
Wall-Clock-Zeiten der drei Varianten gilt $t_1 \ge t_2 \ge t_3$; eine Verletzung ist ein Mess-Defekt
der betroffenen Zelle, kein Ergebnis, und führt zum Ausschluss der Zelle. Durchführungs-Stand: Diese
Plausibilitäts-Prüfung ist im Code nicht gebaut; sie schließt an die Messfehler-Interpolation als
fünfte Auswertungs-Komponente der compare-Stufe an (Abschnitt~\ref{sec:heuristics-curves}), deren
Konstellationen der Planer dynamisch aus der Torlage rechnet, nie als statischer Nenner.

Die zweite Stufe ist der komplementäre \emph{vierte Schritt} der finalen heuristischen
Zusammenstellung: Die dem Heuristik-Tier untergebenen Tier-Binaries --- nicht die final optimierte
Haupt-Binary --- werden ohne Mikro- und Makro-Benchmarks und ohne Beobachter neu gebaut, sodass kein
Mess-Artefakt im Kompilat verbleibt. Verglichen wird dann rein an der Wall-Clock, in zwei
Vergleichsmengen mit festgelegtem Vorrang: primär gegen die rekonstruierten Paper-Algorithmen --- die
33 Korpus-Arbeiten mit je einem Profil-XML, 30 Vollprofile und 3 abstrakte
(Abschnitt~\ref{sec:sota-instances}); der Paper-Vergleich geht vor, denn er beantwortet die
Bewertungsfrage wörtlich ---, sekundär gegen alle bekannten Tier-Binary-Permutationen, insbesondere
das beste Einzel-Binary; dieser zweite Vergleich ordnet ein, ob ein etwaiger Vorsprung der Heuristik
als Ganzer oder schlicht einer einzeln überlegenen Permutation zuzuschreiben ist.

Daraus folgt die Entscheidungsregel: Der Gegenbeweis gilt je Workload-Klasse als \emph{erbracht},
wenn die beobachterfrei neu gebaute Rekombination die schnellste Paper-Rekonstruktion in der
Wall-Clock schlägt und der Abstand das Unsicherheitsband der Bau-Varianten übersteigt; er gilt als
\emph{verfehlt}, wenn ein Paper-Algorithmus oder das beste Einzel-Binary vorn liegt --- auch dieses
Ergebnis wäre ein Resultat, denn es bezifferte erstmals, wie viel Schätzfehler zwischen
Kurven-Prognose und realem Verhalten liegt. Das Unsicherheitsband ist von einem Mess-Defekt zu
unterscheiden: Zellen, die das Gültigkeitskriterium der ersten Stufe verletzen, gehen nicht in die
Entscheidung ein. Die Messdaten dieser Dimension liegen nicht vor; sie werden mit den Reihen A--C in
Abschnitt~\ref{sec:eval-available} ausgewiesen.

\section{Ergebnisse und Analyse}\label{sec:eval-results}
\subsection{Identität der Ergebnisse: Stempel, Lager, Mappe}\label{sec:eval-identity}
%% OWNER-VORLAGE: OV-04 --- Ersatz-Tag für das entfallene e-Suffix: kein neues Versions-Suffix; die
%% Misch-Strategie reist über die merge_mode-Token und die hierarchisch erweiterten Algorithmus-Namen der
%% Organ-Zeile (gleichlautend mit Kapitel 3). Default gesetzt; Widerspruch genügt.
Jede Ergebniszeile ist über das Stempelsystem identifizierbar (Abschnitt~\ref{sec:stamp-model}): die
\texttt{binary\_id} --- der Pfad aus achtzehn \texttt{achse=wert}-Segmenten in Slot-Reihenfolge ---,
der SHA-512-Fingerprint über elf Glieder mit der Format-Kennung \texttt{fingerprint\_format=6} und die
drei einkompilierten Stempel-Zeilen je Realm (Organ, System, Mess). Diese Identität reist unverändert
durch die Kette: als Lager-Schlüssel (Bestand~1 und~2), als Spalten \texttt{permutation\_id} und
\texttt{fingerprint} der Mess-Zeile, als Blatt-Schlüssel der Mappe und als Kennung in jedem Diagramm
und jeder Tabelle des Anhangs --- so bleibt jede dargestellte Zahl auf ihre Binary, ihre
Achsen-Versionen und ihre Maschine zurückführbar. Die Misch-Strategie eines Prüflings wird nicht über
ein Versions-Suffix getragen, sondern über die \texttt{merge}-Token \texttt{replace}, \texttt{merge}
und \texttt{union} der Anwender-XML und die hierarchisch erweiterten Algorithmus-Namen der Organ-Zeile
unter dem Namensraum-Präfix des Prüflings; das Flag \texttt{e} der Eintrags-Grammatik bezeichnet einen
Efficiency-Core-Algorithmus (den E-Kern-Pfad einer Hybrid-CPU), kein Experiment und keine Mischung.
Werte der Messung erscheinen nie im Stempel; gestempelt wird die Übersetzungszeit-Identität.
Durchführungs-Stand: Die Anhang-Tabellen konsumieren \texttt{permutation\_id} und
\texttt{fingerprint} je Zeile; die volle Stempel-Identität je Diagramm ist Ausbaustufe.

\subsection{Datenfluss von der Messung zum Manuskript}\label{sec:eval-dataflow}
Je Reihen-Lauf schreibt die Messstufe die eine offizielle \texttt{measurements.csv} als Projektion der
Mappe; die \texttt{<output>}-Pfade der Anwender-XML (Muster \texttt{\_runs/<Datum>/<Reihe>/}) löst der
Laufeintrag als Lauf-Provenienz auf. Drei Spaltenlisten sind zu unterscheiden
(Tabelle~\ref{tab:eval-schema-ebenen}): das WIDE-Schema mit 189 Spalten (Trenner Semikolon; die
Schema-Hoheit der Mess-Ausgabe, deren Kopf zur Laufzeit aus vier Einzelquellen entsteht und deshalb als
Freeze verankert ist), die Pipeline-Sicht mit 16 Spalten plus 16 Zusatzspalten, also 32, für die
LaTeX-Stufen (Trenner Komma) und die zwölfspaltige Altform des Ergebnis-Aggregators über
\texttt{ExecutionResult}. Die Pipeline-Sicht trägt je Permutation \texttt{permutation\_id},
\texttt{fingerprint}, \texttt{succeeded}, \texttt{workload\_used}, \texttt{op\_count},
\texttt{total\_cycles}, \texttt{cache\_misses\_l1}, \texttt{cache\_misses\_l2},
\texttt{cache\_misses\_l3}, \texttt{dtlb\_misses}, \texttt{coherence\_\allowbreak{}invalidations},
\texttt{energy\_\allowbreak{}micro\_\allowbreak{}joules}, \texttt{bytes\_\allowbreak{}allocated}, \texttt{bytes\_\allowbreak{}in\_\allowbreak{}use\_\allowbreak{}peak},
\texttt{external\_\allowbreak{}frag} und \texttt{internal\_\allowbreak{}frag}; die volle Sicht ergänzt sechs Such-Zähler
(\texttt{search\_\allowbreak{}insert} bis \texttt{search\_\allowbreak{}peak\_\allowbreak{}occupancy}), \texttt{pmc\_\allowbreak{}available},
\texttt{branch\_\allowbreak{}misses}, \texttt{throughput\_\allowbreak{}ops\_\allowbreak{}per\_\allowbreak{}sec} und die sieben
\texttt{*\_source\_available}-Flags. \texttt{total\_cycles} trägt --- unter seinem Legacy-Namen ---
die repräsentative p50-Latenz in Nanosekunden, keine Zyklen; die Umbenennung oder ein echter
Zyklen-Kanal ist ein benannter Code-Entscheid. Nicht erhobene Werte stehen als n/a, nie als Null.
\begin{table}[htbp]\centering\footnotesize
\begin{tabular}{@{}l>{\raggedright\arraybackslash}p{1.5cm}>{\raggedright\arraybackslash}p{1.9cm}>{\raggedright\arraybackslash}p{4.6cm}>{\raggedright\arraybackslash}p{4.0cm}@{}}
\toprule
Schema & Trenner & Spalten & Hoheit / Konsument & Quelle im Code \\
\midrule
WIDE (E4) & Semikolon & 189 & Mess-Ausgabe; Kopf zur Laufzeit aus vier Einzelquellen, als Freeze verankert & Iterator des CacheEngineBuilders (\texttt{lazy\_csv\_header}); Schema-Freeze \\
PIPELINE & Komma & 16 (+16 = 32) & LaTeX-Stufen 04/05/06; volle Sicht als Präfix-Superset & Pipeline-CSV-Schema: eine Liste, vier Schreiber \\
F15-ALT & Komma & 12 & Ergebnis-Aggregator über \texttt{ExecutionResult} & Ergebnis-Aggregator der Builder-Kommandos \\
\bottomrule
\end{tabular}
\caption{Die drei Mess-Spaltenlisten der Cache-Engine und ihre Hoheit. \texttt{total\_cycles} der
Pipeline-Sicht trägt die repräsentative p50-Latenz in Nanosekunden (Legacy-Name), keine
Zyklen.}\label{tab:eval-schema-ebenen}
\end{table}

Die Mappe (xlsx) entsteht immer --- im Speicher, bedingungslos, auch wenn ein Profil nur CSV
persistiert; die CSV ist ihr Kind: Eine Zeile erreicht die CSV nur, wenn die Mappe sie zuvor
angenommen hat, und bricht die Mappen-Erzeugung, ist auch die CSV ungültig; was auf die Platte geht,
entscheidet allein \texttt{<writeback\_methods>} --- xlsx, csv oder beide, als Strategie über den
Ausgabe-Formaten mit xlsx als Default (Abschnitt~\ref{sec:impl-system-axes}); der xlsx-Writer liegt in
der Lager-Ablage und schreibt über libxlsxwriter. Die Auswertung reichert die Projektion über
Binary-zu-CSV um Profil-ID, Paper-Referenz und Achsen-Belegung an, rendert CSV-zu-LaTeX-Steckbriefe
je Permutations-ID (eine leere Eingabe erzeugt keine Datei, nie einen leeren Plot) und
Diagramm-Generator-Diagramme als TikZ-Balken, Scatter-Plots und Heatmaps mit A4-Deckel und
Größen-Wache (eine 1$\times$1-, 1$\times$n- oder n$\times$1-Matrix erzeugt kein Bild); der
Anhang-Generator schreibt je Sprache in ein \texttt{tabellen}-Verzeichnis, aus dem das Manuskript die
Dateien einbindet --- die deutsche Fassung bindet 38 erzeugte Dateien des Anhangs ein.

%% OWNER-VORLAGE: OV-02 --- Mappen-Aufmachung: Writer und Projektion gebaut; Contract-Test an der
%% Standard-xlsx (Owner-Soll) gegen die E1-Prüfung an der CSV (Ist); Export-Element, Chart-Features,
%% dritte Factory LaTeX, Warnungs-Sektion des Info-Blatts und WIDE-Freeze-Nachzug als Ausbaustufen.
%% Durchführungs-Stand an dieser einen Stelle; Widerspruch genügt.
Die Mappe eines Laufs ist eine Datei mit einem Blatt je gewählter Unter-Achsen-Permutation und einem
Info-Blatt (System-Informationen der messenden Maschine, verwendete Haupt-Achsen aller drei Realms,
konstante Meta-Achsen); die Blätter sind in der bindenden Ordnung Mess $\to$ System $\to$ Organ
nummeriert (\texttt{S001} bis \texttt{Snnn}), und die Spalten stammen ausschließlich aus dem
WIDE-Kopf; Spalten, deren Wert sich innerhalb eines Laufs nie ändert, wandern in die Metadaten des
Info-Blatts, statt in jeder Zeile wiederholt zu werden. Die Contract-Tests der Persistenz sollen an der
Standard-xlsx prüfen, nicht an der CSV (Abschnitt~\ref{sec:impl-contracts}). Neben xlsx und csv ist
eine dritte Factory LaTeX vorgesehen, die die Blatt-Ordnung der Mappe beim Ausrollen des Anhangs
serialisiert und je Binary eine Tabellenzeile als zeitliche Messreihe führt. Durchführungs-Stand:
Writer und Projektion sind gebaut; die Naht schreibt ein Datenblatt unter konstantem Blattschlüssel, die
Auffächerung in ein Blatt je Unter-Achsen-Permutation plus Info-Blatt ist Soll
(Abschnitt~\ref{sec:limitations}); die Persistenz-Prüfung liest die CSV-Ausgabe, die Prüfung an der
je Binary erzeugten xlsx ist als Bau-Schritt benannt; das Export-Element mit Ziel-Filter und
per-Binary-xlsx, die Chart-Features, die LaTeX-Factory, die Warnungs-Sektion des Info-Blatts und der
Nachzug des WIDE-Freeze sind Ausbaustufen.

\subsection{Messreihen und Verbund-Strategien}\label{sec:eval-series}
%% OWNER-VORLAGE: OV-04 / OV-05 --- kein e-Suffix für die Misch-Strategie (e = Efficiency Core);
%% Begriffs-Kanon: die merge-Token replace/merge/union führen, Verbund-Namen als Haus-Begriff, Stufe 1..3 nur
%% Alias, Full Join auf union zurückgeführt. Defaults gesetzt; Widerspruch genügt.
Drei Messreihen werden ausgewertet (Tabelle~\ref{tab:stage-series}): Reihe~A --- PRT-ART gegen den
Stand der Technik mit den Evaluations-Hypothesen H1--H3 (Abschnitt~\ref{sec:eval-hypotheses}); Reihe~B
--- die systematische Variation der Cache-Engine-Permutationen, gespeist aus der Union der Stufe~3;
Reihe~C --- Merge/Regression alt gegen neu, bau-übergreifend. Quer zu dieser Zuordnung deklariert die
Anwender-XML je Achse den \texttt{merge}-Modus des Prüflings (Tabelle~\ref{tab:eval-verbund-modi}):
\texttt{replace} (der Standard, auch bei leerem Attribut; die Prüflings-Achse ersetzt die
Cache-Engine-Achse mit Rückfall, Verbund-Strategie \texttt{Verbund2\_Replace}), \texttt{merge}
(Cache-Engine und Prüfling hybrid je Prüfling, \texttt{Verbund2\_Hybrid}) und \texttt{union} (die
explizite, an Stufe~3 gebundene Vereinigung aller Achsen-Varianten ohne Verwerfen,
\texttt{Verbund3\_Union}); Stufe~1 ist \texttt{Verbund1\_CeOnly}, die Abwesenheit jeder
Prüflings-Direktive, und \enquote{Stufe~1 bis~3} bleibt der Alias der Verbund-Namen, der Full Join
der Alias von \texttt{union}. Reihe~A arbeitet mit \texttt{replace} und \texttt{merge}, Reihe~B mit
\texttt{union}. Unterscheidbar bleiben die Stände an der Binary über die hierarchisch erweiterten
Algorithmus-Namen der Organ-Zeile (Abschnitt~\ref{sec:eval-identity}); eine eigene Stempel-Zeile
tragen sie nicht, denn der Prüfling mischt gegen die Organ-Achsen und steht damit bereits in deren
Zeile, und ein Versions-Suffix trägt die Mischung nicht. Bei reinen Ersetzungs- wie reinen
Cache-Engine-Bauten bleibt die Namens-Erweiterung folgerichtig aus. Durchführungs-Stand: Das Token
\texttt{merge} ist im Merge-Plan registriert und im golden-Kern-Profil deklariert; materialisiert
sind die Strategien \texttt{Verbund1\_CeOnly}, \texttt{Verbund2\_Replace} und
\texttt{Verbund3\_Union}, während \texttt{Verbund2\_Hybrid} ohne Produktions-Aufrufer bleibt, bis
seine Materialisierung im Strategie-Enum landet.
\begin{table}[htbp]\centering\footnotesize
\begin{tabular}{@{}ll>{\raggedright\arraybackslash}p{2.2cm}>{\raggedright\arraybackslash}p{4.4cm}>{\raggedright\arraybackslash}p{1.6cm}@{}}
\toprule
\texttt{merge}-Token & Verbund-Strategie & Alias & Wirkung & Reihe \\
\midrule
--- (keine Direktive) & \texttt{Verbund1\_CeOnly} & Stufe~1 & nur Cache-Engine: Standard-Bausteine und SOTA-Profile & A (Baseline, auch C) \\
leer / \texttt{replace} & \texttt{Verbund2\_Replace} & Stufe~2 (Ersetzen) & die Prüflings-Achse ersetzt die Cache-Engine-Achse, mit Rückfall & A \\
\texttt{merge} & \texttt{Verbund2\_Hybrid} & Stufe~2 (Hybrid) & Cache-Engine und Prüfling hybrid je Prüfling; Materialisierung offen, kein Produktions-Aufrufer & A \\
\texttt{union} & \texttt{Verbund3\_Union} & Stufe~3 (Full Join) & Vereinigung aller Achsen-Varianten ohne Verwerfen, an Stufe~3 gebunden & B \\
\bottomrule
\end{tabular}
\caption{Die \texttt{merge}-Token der Anwender-XML und ihre Verbund-Strategien im Merge-Plan; die
Stufen-Namen sind Aliasse, der Full Join ist der Alias von \texttt{union}. Die Reihen-Spalte folgt
der Zuordnung aus Tabelle~\ref{tab:stage-series}.}\label{tab:eval-verbund-modi}
\end{table}

\subsection{Vorliegende Messungen}\label{sec:eval-available}
Methodik und Auswertungs-Pipeline sind vorhanden und an realen Daten verifiziert: Eine Smoke-Messreihe
über 43 Permutationen --- echte DLL-Messung; 27 kartesische Kombinationen aus SIMD, Speicherlayout
und Allokator plus 16 weitere über vier zweiwertige Größen --- dokumentiert
Anhang~\ref{app:measurements} samt ihren Limitierungen
(Abschnitt~\ref{sec:measurements:limitations}); sie variiert sieben Größen (SIMD ist eine
System-Achse, Speicherlayout und Allokator sind Organ-Achsen, die übrigen vier Namen sind keine
Registry-Bausteine und werden nicht gleichgesetzt), die Kampagne variiert alle achtzehn
Kompositions-Achsen. Die vollständigen Läufe der Pflicht-Messreihen A--C auf den
Zielplattformen liegen nicht vor: n/a, weil der Voll-Bau --- der Trigger der Kampagne --- nicht
gefahren ist (Abschnitt~\ref{sec:limitations}); dasselbe gilt für die Messdaten des Gegenbeweises
(Abschnitt~\ref{sec:eval-counterproof}). Keine Zahl dieses Kapitels stammt aus einem anderen Lauf als
dem im Anhang ausgewiesenen.

% Einbau nach dem ersten vollständigen Lauf der Reihen A--C (Generator-Ausgabe je Sprache):
%   \input{anhang/de/tabellen/<name>.tex} -- Reihe A_full (belegt: die 33 Experiment-XML deklarieren A_full);
%   die Dateinamen der Reihen B und C entstehen mit dem Lauf.

\subsection{Achsen-Sensitivitätsanalyse}\label{sec:eval-sensitivity}
%% OWNER-VORLAGE: OV-K5-02 --- die vierte Dimension (Gegenbeweis) steht neben der Rangbildung über die drei
%% Mess-Ebenen, nicht in ihr. Default gesetzt; Widerspruch genügt.
Den Kern der Auswertung bildet die \emph{Achsen-Sensitivitätsanalyse}: Sie beantwortet das
Trennbarkeits-Problem der Aufgabenstellung (Abschnitte~\ref{sec:rqs} und~\ref{sec:task-concerns}),
indem das Achsen-Framework jeden Bestandteil als austauschbare Achse über einer sonst konstanten
Struktur isoliert und der Beitrag jeder Entwurfsentscheidung unter sonst gleichen Bedingungen
zurechenbar wird. Die Rangbildung läuft über die drei Mess-Ebenen in Kanon-Reihenfolge:
\emph{wallclock} --- das Gesamt-Benchmarking der YCSB-Lastprofile im CacheEngineBuilder ---,
\emph{macro} --- die Genus-Interface-Operationen der Tier-Binary --- und \emph{micro} --- das
Achsen-Interface, die Funktionsmenge, die alle Algorithmen einer Achse liefern müssen
(Abschnitt~\ref{sec:mess-chain}); anhand der Zeilen je Permutation quantifiziert die Analyse auf der
micro-Ebene, welche Achse am meisten zur Performance-Varianz beiträgt, die wallclock-Ebene prüft die
interne micro-Messung von außen, und die Mess-Kategorien sind dabei die Unter-Achse der Tooling-Achse
(Abschnitt~\ref{sec:eval-counters}). Leitachse ist die Cacheline-Awareness: welche Achse die
cache-nahen Kategorien auf welcher Plattform bewegt. Aus der Rangbildung folgt die Empfehlung von
Standard-Konfigurationen je Workload-Klasse; wo Verhältnisse über Workloads oder Plattformen
zusammengefasst werden, geschieht das geometrisch, nie arithmetisch
gemittelt~\cite{fleming1986statistics}. Der heuristische Gegenbeweis steht neben dieser Rangbildung
(Abschnitt~\ref{sec:eval-counterproof}).

\subsection{Gültigkeit und Grenzen}\label{sec:eval-validity}
Die Gültigkeit der Evaluation wird nach den vier Klassen der empirischen Softwareforschung
geprüft\footnote{Threats to Validity: die in der empirischen Softwareforschung übliche Einteilung der
Bedrohungen eines Schlusses in Konstrukt-Validität (misst die Messgröße das Gemeinte?), interne
Validität (ist der Effekt kausal isoliert?), externe Validität (überträgt sich der Befund auf andere
Plattformen und Lasten?) und Schluss-Validität (tragen Statistik und Aggregation den
Schluss?)~\cite{wohlin2012experimentation}.} (Tabelle~\ref{tab:eval-gueltigkeit};
Abschnitt~\ref{sec:measurement-system} führt dieselben Klassen als Gütekriterien des Mess-Systems).
Konstrukt-Validität: Messen die Kategorien die Cacheline-Awareness? Die PMC-Semantik ist je Modell
offengelegt (Abschnitt~\ref{sec:eval-counters}), und Hardware-Zähler gelten nur kreuzgeprobt als
vertrauenswürdig~\cite{weaver2008counters}. Interne Validität: Der Beobachter-Overhead ist über das
Unsicherheitsband der drei Bau-Varianten gebunden (Abschnitt~\ref{sec:eval-counterproof}); Layout- und
Umgebungs-Bias~\cite{mytkowicz2009wrongdata,curtsinger2013stabilizer} begegnen identische Schlüssel,
identische Bau-Matrix und Wiederholungen, eine Layout-Randomisierung fehlt; Warmup und
Wiederholungen folgen Abschnitt~\ref{sec:fairness}~\cite{kalibera2013rigorous}. Externe Validität: Die
registrierte Flotte umfasst zwei x86-64-Mikroarchitekturen und einen Gracemont-Knoten; ZIH und HBM
sind deklariert, nicht vermessen (Abschnitt~\ref{sec:eval-platforms}), die Intel-RAW-Belegung fehlt,
und eine einseitige Lane-Fahrt gilt als Teil-Messung. Schluss-Validität: Perzentile nearest-rank ohne
Interpolation, geometrische Aggregation, n/a statt Null. Die Limitierungen des Mess-Anhangs
(Tabelle~\ref{tab:le:limitierung}, Abschnitt~\ref{sec:measurements:limitations}) und des Fazits
(Abschnitt~\ref{sec:limitations}) sind die belegte Liste dieser Grenzen; keine Zahl wird ohne Lauf
genannt, und die Berichtsregeln für Leistungsmessungen~\cite{hoefler2015benchmarking,blackburn2016truth}
gelten für jede Zahl, die ein Lauf liefert.
\begin{table}[htbp]\centering\footnotesize
\begin{tabular}{@{}>{\raggedright\arraybackslash}p{1.6cm}>{\raggedright\arraybackslash}p{4.0cm}>{\raggedright\arraybackslash}p{4.9cm}>{\raggedright\arraybackslash}p{3.7cm}@{}}
\toprule
Klasse & Bedrohung & Gegenmaßnahme im Apparat & Stand \\
\midrule
Konstrukt & Zähler messen nicht das Gemeinte (Modell-Semantik, Prefetch-Anteile) & RAW-Katalog je Modell mit Kreuzprobe; Semantik-Offenlegung je Spalte; n/a ohne Katalog-Eintrag & gebaut (Zen~5); Intel-Eintrag: Ausbaustufe \\
Konstrukt & Zeit-Kategorie als Zyklen fehlgelesen & \texttt{total\_cycles} als p50-Latenz in Nanosekunden deklariert & gebaut; Umbenennung: Code-Entscheid \\
Intern & Beobachter-Overhead verfälscht die Kurven & drei Bau-Varianten, Unsicherheitsband, Gültigkeitskriterium $t_1 \ge t_2 \ge t_3$ & Gate-Felder gebaut; Prüfung: Ausbaustufe \\
Intern & Layout-/Umgebungs-Bias, Wechsel des L3-Komplexes & identische Schlüssel und Bau-Matrix; Wiederholungs-Dimension; Pinning als Soll deklariert & Aktuator gebaut, Verdrahtung: Ausbaustufe; keine Layout-Randomisierung \\
Intern & Warmup- und Kaltstart-Effekte & Zwei-Phasen-Treiber (sichern, aufwärmen, zurückrollen, messen) & gebaut \\
Extern & Flotte deckt die Zielplattformen nicht ab & registrierte Signaturen prod1, prod2, ODROID\@; ZIH als IS3 deklariert & ZIH und HBM\@: deklariert, nicht vermessen (n/a) \\
Extern & einseitige Hersteller-Lane & Vollständigkeits-Regel der Auswertung; Teil-Messung wird ausgewiesen & Regel gesetzt; Intel-Lane: Kreuzprobe ausstehend \\
Schluss & Perzentile erfunden, Mittelwerte verzerrt & Nearest-Rank ohne Interpolation; geometrische Aggregation & gebaut (Auswertungs-Baum); Kopien in den Werkzeugen: benannte Lücke \\
Schluss & Null statt Nichterhebung & Feld-Flags je Zähler, n/a Ende-zu-Ende & gebaut \\
\bottomrule
\end{tabular}
\caption{Gültigkeit der Evaluation nach den vier Klassen der Bedrohungen (Konstrukt, intern, extern,
Schluss): je Bedrohung die Gegenmaßnahme im Apparat und ihr Durchführungs-Stand --- gebaut,
deklariert und nicht vermessen (n/a) oder Ausbaustufe. Keine Zeile nennt eine Zahl ohne
Lauf.}\label{tab:eval-gueltigkeit}
\end{table}

\subsection{Zusammenfassung der Evaluation}\label{sec:eval-summary}
Die Evaluation hält je Evaluationsobjekt fest. Entwurfsraum: Die Hypothesen H1--H3 tragen Messgröße und
Kriterium, und die Sensitivitätsanalyse ist als Verfahren über die drei Mess-Ebenen festgelegt; Belege
liegen für die Smoke-Reihe vor (43 Permutationen, sieben Größen, Anhang~\ref{app:measurements}), für die
Pflicht-Messreihen nicht (n/a, Voll-Bau ausstehend). Apparat: Der Gegenbeweis ist als vierte Dimension
mit Prüfplan, Gültigkeitskriterium und Entscheidungsregel festgelegt, seine Messdaten liegen nicht vor;
das Messfehler-Fundament --- Gate-Felder im Fingerprint, Feld-Flags je Zähler, n/a-Regel --- ist
gebaut. Automatisierung: Die Kette von der Experiment-XML bis zum eingebundenen Anhang ist am
Durchstich real gelaufen und liefert jede Zahl mit ihrem Lauf. Die Ergebnisse werden vor dem
Leitmotiv der Arbeit diskutiert --- dem Übergang vom heuristischen zum informierten Cache-Management,
also der Frage, ob telemetrie-getriebene, automatisch adaptierte Konfigurationen statisch gewählte
übertreffen; der Cache-Line-Effekt aus den Synthese-Kurven ist der erste Auswertungsschritt nach dem
vollständigen Lauf, datenabhängig und deshalb hier nicht vorweggenommen. Nach den Maßstäben für
Forschungsbeiträge~\cite{shaw2003writing} und für empirische Bewertungen~\cite{blackburn2016truth} ist
der Beitrag dieses Kapitels ein validierter Versuchsaufbau samt seiner Automatisierung, dessen
Ergebnis-Zahlen der Kampagne folgen. Kapitel~\ref{ch:conclusion} beantwortet die Forschungsfragen auf
dieser Grundlage.
